
Un chariot sur roues de \qty{\v{Masse}}{g} (frottement négligeable) oscille à l’extrémité d’un ressort, sur une surface horizontale en effectuant un mouvement harmonique simple. La constante de rappel du ressort est de \qty{\v{K}}{N/m}. À $t = 0\text{ s}$, le chariot, déjà en mouvement, se trouve à \qty{\v{PosInit}}{cm} à droite de sa position d’équilibre et sa vitesse est de \qty{\v{VitInit}}{cm/s} vers la droite.

\begin{enumerate}[label=\alph*)]
    \item Déterminez l’équation de sa position en fonction du temps $x(t)$ à l’aide d’une fonction sinus. Justifiez la valeur trouvée pour la constante de phase.
    \item À quel(s) moment(s) dans un cycle l’énergie potentielle du chariot est-elle maximale ? Expliquez votre réponse qualitativement, sans aucun calcul.
\end{enumerate}